\documentclass[
  doc,
  floatsintext,
  longtable,
  a4paper,
  nolmodern,
  notxfonts,
  notimes,
  12pt,
  colorlinks=true,linkcolor=blue,citecolor=blue,urlcolor=blue]{apa7}

\usepackage{amsmath}
\usepackage{amssymb}

\geometry{inner=1in, outer=1in}
\fancyhfoffset[LE,RO]{0cm}


\usepackage[bidi=default]{babel}
\babelprovide[main,import]{spanish}


% get rid of language-specific shorthands (see #6817):
\let\LanguageShortHands\languageshorthands
\def\languageshorthands#1{}

\RequirePackage{longtable}
\RequirePackage{threeparttablex}

\makeatletter
\renewcommand{\paragraph}{\@startsection{paragraph}{4}{\parindent}%
	{0\baselineskip \@plus 0.2ex \@minus 0.2ex}%
	{-.5em}%
	{\normalfont\normalsize\bfseries\typesectitle}}

\renewcommand{\subparagraph}[1]{\@startsection{subparagraph}{5}{0.5em}%
	{0\baselineskip \@plus 0.2ex \@minus 0.2ex}%
	{-\z@\relax}%
	{\normalfont\normalsize\bfseries\itshape\hspace{\parindent}{#1}\textit{\addperi}}{\relax}}
\makeatother




\usepackage{longtable, booktabs, multirow, multicol, colortbl, hhline, caption, array, float, xpatch}
\usepackage{subcaption}


\renewcommand\thesubfigure{\Alph{subfigure}}
\setcounter{topnumber}{2}
\setcounter{bottomnumber}{2}
\setcounter{totalnumber}{4}
\renewcommand{\topfraction}{0.85}
\renewcommand{\bottomfraction}{0.85}
\renewcommand{\textfraction}{0.15}
\renewcommand{\floatpagefraction}{0.7}

\usepackage{tcolorbox}
\tcbuselibrary{listings,theorems, breakable, skins}
\usepackage{fontawesome5}

\definecolor{quarto-callout-color}{HTML}{909090}
\definecolor{quarto-callout-note-color}{HTML}{0758E5}
\definecolor{quarto-callout-important-color}{HTML}{CC1914}
\definecolor{quarto-callout-warning-color}{HTML}{EB9113}
\definecolor{quarto-callout-tip-color}{HTML}{00A047}
\definecolor{quarto-callout-caution-color}{HTML}{FC5300}
\definecolor{quarto-callout-color-frame}{HTML}{ACACAC}
\definecolor{quarto-callout-note-color-frame}{HTML}{4582EC}
\definecolor{quarto-callout-important-color-frame}{HTML}{D9534F}
\definecolor{quarto-callout-warning-color-frame}{HTML}{F0AD4E}
\definecolor{quarto-callout-tip-color-frame}{HTML}{02B875}
\definecolor{quarto-callout-caution-color-frame}{HTML}{FD7E14}

%\newlength\Oldarrayrulewidth
%\newlength\Oldtabcolsep


\usepackage{hyperref}



\usepackage{color}
\usepackage{fancyvrb}
\newcommand{\VerbBar}{|}
\newcommand{\VERB}{\Verb[commandchars=\\\{\}]}
\DefineVerbatimEnvironment{Highlighting}{Verbatim}{commandchars=\\\{\}}
% Add ',fontsize=\small' for more characters per line
\usepackage{framed}
\definecolor{shadecolor}{RGB}{241,243,245}
\newenvironment{Shaded}{\begin{snugshade}}{\end{snugshade}}
\newcommand{\AlertTok}[1]{\textcolor[rgb]{0.68,0.00,0.00}{#1}}
\newcommand{\AnnotationTok}[1]{\textcolor[rgb]{0.37,0.37,0.37}{#1}}
\newcommand{\AttributeTok}[1]{\textcolor[rgb]{0.40,0.45,0.13}{#1}}
\newcommand{\BaseNTok}[1]{\textcolor[rgb]{0.68,0.00,0.00}{#1}}
\newcommand{\BuiltInTok}[1]{\textcolor[rgb]{0.00,0.23,0.31}{#1}}
\newcommand{\CharTok}[1]{\textcolor[rgb]{0.13,0.47,0.30}{#1}}
\newcommand{\CommentTok}[1]{\textcolor[rgb]{0.37,0.37,0.37}{#1}}
\newcommand{\CommentVarTok}[1]{\textcolor[rgb]{0.37,0.37,0.37}{\textit{#1}}}
\newcommand{\ConstantTok}[1]{\textcolor[rgb]{0.56,0.35,0.01}{#1}}
\newcommand{\ControlFlowTok}[1]{\textcolor[rgb]{0.00,0.23,0.31}{\textbf{#1}}}
\newcommand{\DataTypeTok}[1]{\textcolor[rgb]{0.68,0.00,0.00}{#1}}
\newcommand{\DecValTok}[1]{\textcolor[rgb]{0.68,0.00,0.00}{#1}}
\newcommand{\DocumentationTok}[1]{\textcolor[rgb]{0.37,0.37,0.37}{\textit{#1}}}
\newcommand{\ErrorTok}[1]{\textcolor[rgb]{0.68,0.00,0.00}{#1}}
\newcommand{\ExtensionTok}[1]{\textcolor[rgb]{0.00,0.23,0.31}{#1}}
\newcommand{\FloatTok}[1]{\textcolor[rgb]{0.68,0.00,0.00}{#1}}
\newcommand{\FunctionTok}[1]{\textcolor[rgb]{0.28,0.35,0.67}{#1}}
\newcommand{\ImportTok}[1]{\textcolor[rgb]{0.00,0.46,0.62}{#1}}
\newcommand{\InformationTok}[1]{\textcolor[rgb]{0.37,0.37,0.37}{#1}}
\newcommand{\KeywordTok}[1]{\textcolor[rgb]{0.00,0.23,0.31}{\textbf{#1}}}
\newcommand{\NormalTok}[1]{\textcolor[rgb]{0.00,0.23,0.31}{#1}}
\newcommand{\OperatorTok}[1]{\textcolor[rgb]{0.37,0.37,0.37}{#1}}
\newcommand{\OtherTok}[1]{\textcolor[rgb]{0.00,0.23,0.31}{#1}}
\newcommand{\PreprocessorTok}[1]{\textcolor[rgb]{0.68,0.00,0.00}{#1}}
\newcommand{\RegionMarkerTok}[1]{\textcolor[rgb]{0.00,0.23,0.31}{#1}}
\newcommand{\SpecialCharTok}[1]{\textcolor[rgb]{0.37,0.37,0.37}{#1}}
\newcommand{\SpecialStringTok}[1]{\textcolor[rgb]{0.13,0.47,0.30}{#1}}
\newcommand{\StringTok}[1]{\textcolor[rgb]{0.13,0.47,0.30}{#1}}
\newcommand{\VariableTok}[1]{\textcolor[rgb]{0.07,0.07,0.07}{#1}}
\newcommand{\VerbatimStringTok}[1]{\textcolor[rgb]{0.13,0.47,0.30}{#1}}
\newcommand{\WarningTok}[1]{\textcolor[rgb]{0.37,0.37,0.37}{\textit{#1}}}

\providecommand{\tightlist}{%
  \setlength{\itemsep}{0pt}\setlength{\parskip}{0pt}}
\usepackage{longtable,booktabs,array}
\newcounter{none} % for unnumbered tables
\usepackage{calc} % for calculating minipage widths
% Correct order of tables after \paragraph or \subparagraph
\usepackage{etoolbox}
\makeatletter
\patchcmd\longtable{\par}{\if@noskipsec\mbox{}\fi\par}{}{}
\makeatother
% Allow footnotes in longtable head/foot
\IfFileExists{footnotehyper.sty}{\usepackage{footnotehyper}}{\usepackage{footnote}}
\makesavenoteenv{longtable}

\usepackage{graphicx}
\makeatletter
\newsavebox\pandoc@box
\newcommand*\pandocbounded[1]{% scales image to fit in text height/width
  \sbox\pandoc@box{#1}%
  \Gscale@div\@tempa{\textheight}{\dimexpr\ht\pandoc@box+\dp\pandoc@box\relax}%
  \Gscale@div\@tempb{\linewidth}{\wd\pandoc@box}%
  \ifdim\@tempb\p@<\@tempa\p@\let\@tempa\@tempb\fi% select the smaller of both
  \ifdim\@tempa\p@<\p@\scalebox{\@tempa}{\usebox\pandoc@box}%
  \else\usebox{\pandoc@box}%
  \fi%
}
% Set default figure placement to htbp
\def\fps@figure{htbp}
\makeatother







\usepackage{newtx}

\defaultfontfeatures{Scale=MatchLowercase}
\defaultfontfeatures[\rmfamily]{Ligatures=TeX,Scale=1}





\title{Comentarios}


\shorttitle{COMENTARIOS}


\usepackage{etoolbox}



\ccoppy{\textcopyright~2026}



\author{Edison Achalma}



\affiliation{
{Escuela Profesional de Economía, Universidad Nacional de San Cristóbal
de Huamanga}}




\leftheader{Achalma}

\date{2026-10-20}


\abstract{Este abstract será actualizado una vez que se complete el
contenido final del artículo. }



\authornote{\par{\addORCIDlink{Edison Achalma}{0000-0001-6996-3364}} 
\par{ }
\par{   El autor no tiene conflictos de interés que revelar.    Los
roles de autor se clasificaron utilizando la taxonomía de roles de
colaborador (CRediT; https://credit.niso.org/) de la siguiente
manera:  Edison Achalma:   conceptualización, metodología, análisis
formal, investigación, recursos, curación de
datos, redacción, visualización, supervisión, administración del
proyecto}
\par{La correspondencia relativa a este artículo debe dirigirse a Edison
Achalma, Escuela Profesional de Economía, Universidad Nacional de San
Cristóbal de Huamanga, Portal Independencia N°
57, Ayacucho, AYA 5001, Perú, Email: \href{mailto:elmer.achalma.09@unsch.edu.pe}{elmer.achalma.09@unsch.edu.pe}}
}

\makeatletter
\let\endoldlt\endlongtable
\def\endlongtable{
\hline
\endoldlt
}
\makeatother

\urlstyle{same}



\hypersetup{
  pdftitle={Comentarios},
  pdfauthor={Autor}
}
\makeatletter
\@ifpackageloaded{caption}{}{\usepackage{caption}}
\AtBeginDocument{%
\ifdefined\contentsname
  \renewcommand*\contentsname{Tabla de contenidos}
\else
  \newcommand\contentsname{Tabla de contenidos}
\fi
\ifdefined\listfigurename
  \renewcommand*\listfigurename{Lista de Figuras}
\else
  \newcommand\listfigurename{Lista de Figuras}
\fi
\ifdefined\listtablename
  \renewcommand*\listtablename{Lista de Tablas}
\else
  \newcommand\listtablename{Lista de Tablas}
\fi
\ifdefined\figurename
  \renewcommand*\figurename{Figura}
\else
  \newcommand\figurename{Figura}
\fi
\ifdefined\tablename
  \renewcommand*\tablename{Tabla}
\else
  \newcommand\tablename{Tabla}
\fi
}
\@ifpackageloaded{float}{}{\usepackage{float}}
\floatstyle{ruled}
\@ifundefined{c@chapter}{\newfloat{codelisting}{h}{lop}}{\newfloat{codelisting}{h}{lop}[chapter]}
\floatname{codelisting}{Listado}
\newcommand*\listoflistings{\listof{codelisting}{Lista de Listados}}
\makeatother
\makeatletter
\makeatother
\makeatletter
\@ifpackageloaded{caption}{}{\usepackage{caption}}
\@ifpackageloaded{subcaption}{}{\usepackage{subcaption}}
\makeatother
\makeatletter
\@ifpackageloaded{fontawesome5}{}{\usepackage{fontawesome5}}
\makeatother

% From https://tex.stackexchange.com/a/645996/211326
%%% apa7 doesn't want to add appendix section titles in the toc
%%% let's make it do it
\makeatletter
\xpatchcmd{\appendix}
  {\par}
  {\addcontentsline{toc}{section}{\@currentlabelname}\par}
  {}{}
\makeatother

%% Disable longtable counter
%% https://tex.stackexchange.com/a/248395/211326

\usepackage{etoolbox}

\makeatletter
\patchcmd{\LT@caption}
  {\bgroup}
  {\bgroup\global\LTpatch@captiontrue}
  {}{}
\patchcmd{\longtable}
  {\par}
  {\par\global\LTpatch@captionfalse}
  {}{}
\apptocmd{\endlongtable}
  {\ifLTpatch@caption\else\addtocounter{table}{-1}\fi}
  {}{}
\newif\ifLTpatch@caption
\makeatother

\begin{document}

\maketitle


\hypertarget{toc}{}
\tableofcontents
\newpage
\section[Introduction]{Comentarios}

\setcounter{secnumdepth}{3}

\setlength\LTleft{0pt}




\subsection{🎯 Motivación}\label{motivaciuxf3n}

El código se escribe una vez, pero se lee muchas veces. Los comentarios
son la forma en que te comunicas con otros programadores (y contigo
mismo en el futuro). Un buen comentario explica el ``por qué'' detrás
del ``qué'', haciendo que tu código sea más fácil de entender, mantener
y depurar. Además, los comentarios son una herramienta poderosa para
probar código temporalmente sin eliminarlo.

\subsection{💡 Intuición}\label{intuiciuxf3n}

Imagina que estás escribiendo un libro de cocina. Las instrucciones (el
código) son los pasos de la receta. Los comentarios son como notas al
margen: ``este paso es importante porque evita que se queme'', ``aquí
puedes sustituir el ingrediente si no lo tienes''. El chef (Python)
ignora completamente tus notas al margen y solo sigue las instrucciones
principales.

\subsection{🧠 Modelo mental}\label{modelo-mental}

Los comentarios en Python:

\begin{enumerate}
\def\labelenumi{\arabic{enumi}.}
\tightlist
\item
  \textbf{Comienzan con \texttt{\#}}: todo lo que sigue en esa línea es
  ignorado por Python.
\item
  \textbf{Pueden estar al inicio de la línea} o \textbf{al final} de una
  línea de código.
\item
  \textbf{No tienen efecto en la ejecución}: son solo para humanos.
\item
  \textbf{No existe sintaxis oficial para comentarios de múltiples
  líneas}, pero hay formas de simularlos.
\end{enumerate}

\subsection{📐 Desarrollo}\label{desarrollo}

\subsubsection{1. Comentario de una
línea}\label{comentario-de-una-luxednea}

El comentario comienza con \texttt{\#} y Python ignora todo lo que sigue
hasta el final de la línea:

\begin{Shaded}
\begin{Highlighting}[]
\CommentTok{\# Esto es un comentario}
\BuiltInTok{print}\NormalTok{(}\StringTok{"Hola mundo"}\NormalTok{)  }\CommentTok{\# Esto también es un comentario}
\end{Highlighting}
\end{Shaded}

\subsubsection{2. Comentarios al final de
línea}\label{comentarios-al-final-de-luxednea}

Útiles para explicar una línea específica:

\begin{Shaded}
\begin{Highlighting}[]
\NormalTok{precio }\OperatorTok{=} \DecValTok{100}      \CommentTok{\# Precio base sin impuestos}
\NormalTok{iva }\OperatorTok{=} \FloatTok{0.21}        \CommentTok{\# IVA del 21\%}
\NormalTok{total }\OperatorTok{=}\NormalTok{ precio }\OperatorTok{*}\NormalTok{ (}\DecValTok{1} \OperatorTok{+}\NormalTok{ iva)  }\CommentTok{\# Cálculo del precio final}
\end{Highlighting}
\end{Shaded}

\subsubsection{3. Desactivar código con
comentarios}\label{desactivar-cuxf3digo-con-comentarios}

Comentar una línea evita que se ejecute:

\begin{Shaded}
\begin{Highlighting}[]
\CommentTok{\# print("Esta línea no se ejecuta")}
\BuiltInTok{print}\NormalTok{(}\StringTok{"Esta línea sí se ejecuta"}\NormalTok{)}
\CommentTok{\# Salida: Esta línea sí se ejecuta}
\end{Highlighting}
\end{Shaded}

Útil para depuración: puedes desactivar temporalmente partes del código.

\subsubsection{4. Comentarios de múltiples
líneas}\label{comentarios-de-muxfaltiples-luxedneas}

Python no tiene un símbolo especial para comentarios de varias líneas,
pero puedes:

\textbf{Opción 1: Usar \texttt{\#} en cada línea}

\begin{Shaded}
\begin{Highlighting}[]
\CommentTok{\# Esto es un comentario}
\CommentTok{\# que ocupa varias líneas}
\CommentTok{\# y explica algo complejo}
\BuiltInTok{print}\NormalTok{(}\StringTok{"Hola"}\NormalTok{)}
\end{Highlighting}
\end{Shaded}

\textbf{Opción 2: Usar cadenas de texto con triple comilla (no
asignadas)}

\begin{Shaded}
\begin{Highlighting}[]
\CommentTok{"""}
\CommentTok{Este es un comentario}
\CommentTok{que ocupa varias líneas}
\CommentTok{Python lo ignora porque no está asignado a una variable}
\CommentTok{"""}
\BuiltInTok{print}\NormalTok{(}\StringTok{"Hola"}\NormalTok{)}
\end{Highlighting}
\end{Shaded}

⚠️ \textbf{Importante:} técnicamente esto es una cadena de texto
(string), no un comentario. Aunque Python la ignora, la buena práctica
es usar \texttt{\#} para comentarios oficiales.

\subsubsection{5. Comentarios en bloque y
documentación}\label{comentarios-en-bloque-y-documentaciuxf3n}

\begin{itemize}
\tightlist
\item
  \textbf{Comentarios de bloque}: explican secciones enteras de código.
\item
  \textbf{Docstrings}: comentarios especiales para documentar funciones
  y clases (usando \texttt{"""} y asignados a la función).
\end{itemize}

\begin{Shaded}
\begin{Highlighting}[]
\KeywordTok{def}\NormalTok{ calcular\_area(base, altura):}
    \CommentTok{"""}
\CommentTok{    Calcula el área de un rectángulo.}

\CommentTok{    Args:}
\CommentTok{        base (float): La base del rectángulo}
\CommentTok{        altura (float): La altura del rectángulo}

\CommentTok{    Returns:}
\CommentTok{        float: El área calculada}
\CommentTok{    """}
    \ControlFlowTok{return}\NormalTok{ base }\OperatorTok{*}\NormalTok{ altura}
\end{Highlighting}
\end{Shaded}

\subsection{🔬 Ejemplos}\label{ejemplos}

\textbf{Mínimo ---} Comentarios simples:

\begin{Shaded}
\begin{Highlighting}[]
\CommentTok{\# Saludo inicial}
\BuiltInTok{print}\NormalTok{(}\StringTok{"Hola"}\NormalTok{)  }\CommentTok{\# Muestra un saludo}
\end{Highlighting}
\end{Shaded}

\textbf{Progresivo ---} Comentarios para explicar lógica:

\begin{Shaded}
\begin{Highlighting}[]
\CommentTok{\# Cálculo de descuento}
\NormalTok{precio\_original }\OperatorTok{=} \DecValTok{120}
\NormalTok{descuento }\OperatorTok{=} \FloatTok{0.15}  \CommentTok{\# 15\% de descuento}

\CommentTok{\# El precio final es el original menos el descuento}
\NormalTok{precio\_final }\OperatorTok{=}\NormalTok{ precio\_original }\OperatorTok{{-}}\NormalTok{ (precio\_original }\OperatorTok{*}\NormalTok{ descuento)}
\BuiltInTok{print}\NormalTok{(precio\_final)  }\CommentTok{\# Muestra: 102.0}
\end{Highlighting}
\end{Shaded}

\textbf{Real ---} Documentación en un script de configuración:

\begin{Shaded}
\begin{Highlighting}[]
\CommentTok{\# ==========================================}
\CommentTok{\# CONFIGURACIÓN DEL SISTEMA}
\CommentTok{\# ==========================================}

\CommentTok{\# Configuración de la base de datos}
\CommentTok{\# Usamos SQLite para desarrollo local}
\NormalTok{DB\_HOST }\OperatorTok{=} \StringTok{"localhost"}
\NormalTok{DB\_PORT }\OperatorTok{=} \DecValTok{5432}

\CommentTok{\# Límites del sistema}
\NormalTok{MAX\_USERS }\OperatorTok{=} \DecValTok{100}  \CommentTok{\# Máximo de usuarios simultáneos}
\NormalTok{TIMEOUT }\OperatorTok{=} \DecValTok{30}     \CommentTok{\# Tiempo de espera en segundos}

\CommentTok{\# DEBUG = True   \# Descomentar para modo depuración}
\NormalTok{DEBUG }\OperatorTok{=} \VariableTok{False}    \CommentTok{\# Modo producción}
\end{Highlighting}
\end{Shaded}

\subsection{⚠️ Errores comunes}\label{errores-comunes}

{\def\LTcaptype{none} % do not increment counter
\begin{longtable}[]{@{}
  >{\raggedright\arraybackslash}p{(\linewidth - 4\tabcolsep) * \real{0.2288}}
  >{\raggedright\arraybackslash}p{(\linewidth - 4\tabcolsep) * \real{0.2941}}
  >{\raggedright\arraybackslash}p{(\linewidth - 4\tabcolsep) * \real{0.4771}}@{}}
\toprule\noalign{}
\begin{minipage}[b]{\linewidth}\raggedright
Error
\end{minipage} & \begin{minipage}[b]{\linewidth}\raggedright
Síntoma
\end{minipage} & \begin{minipage}[b]{\linewidth}\raggedright
Cómo detectarlo
\end{minipage} \\
\midrule\noalign{}
\endhead
\bottomrule\noalign{}
\endlastfoot
\textbf{Olvidar \texttt{\#}} & El texto se interpreta como código y da
error & Asegúrate de que todo comentario comienza con \texttt{\#} \\
\textbf{Espacio después de \texttt{\#}} & No da error, pero es una
convención & Pon un espacio después de \texttt{\#} para legibilidad:
\texttt{\#\ Esto\ es\ un\ comentario} \\
\textbf{Confundir comentario con string} &
\texttt{"Esto\ no\ es\ un\ comentario"} es código válido & Los
comentarios usan \texttt{\#}, no comillas \\
\textbf{Anidar comentarios} & No se puede anidar \texttt{\#} dentro de
\texttt{\#} & Cada línea es independiente \\
\end{longtable}
}

\subsection{✅ Buenas prácticas}\label{buenas-pruxe1cticas}

\begin{itemize}
\item
  \textbf{Comenta el ``por qué'', no el ``qué''}: el código ya dice qué
  hace.

\begin{Shaded}
\begin{Highlighting}[]
\CommentTok{\# ❌ Poco útil}
\NormalTok{x }\OperatorTok{=}\NormalTok{ x }\OperatorTok{+} \DecValTok{1}  \CommentTok{\# Incrementa x en 1}

\CommentTok{\# ✅ Útil}
\NormalTok{x }\OperatorTok{=}\NormalTok{ x }\OperatorTok{+} \DecValTok{1}  \CommentTok{\# Contador de intentos fallidos}
\end{Highlighting}
\end{Shaded}
\item
  \textbf{Mantén los comentarios actualizados}: un comentario
  desactualizado es peor que ninguno.
\item
  \textbf{Usa comentarios de bloque} para secciones largas:

\begin{Shaded}
\begin{Highlighting}[]
\CommentTok{\# ==========================================}
\CommentTok{\# SECCIÓN: VALIDACIÓN DE DATOS}
\CommentTok{\# ==========================================}
\end{Highlighting}
\end{Shaded}
\item
  \textbf{Documenta funciones con docstrings} (triple comilla):

\begin{Shaded}
\begin{Highlighting}[]
\KeywordTok{def}\NormalTok{ saludar(nombre):}
    \CommentTok{"""Saluda al usuario por su nombre."""}
    \BuiltInTok{print}\NormalTok{(}\SpecialStringTok{f"Hola, }\SpecialCharTok{\{}\NormalTok{nombre}\SpecialCharTok{\}}\SpecialStringTok{"}\NormalTok{)}
\end{Highlighting}
\end{Shaded}
\item
  \textbf{No abuses de los comentarios}: el código limpio y
  autoexplicativo es mejor que el código lleno de comentarios.
\end{itemize}

\subsection{📝 Ejercicio}\label{ejercicio}

\textbf{Básico} --- Agrega comentarios apropiados

El siguiente código funciona pero no tiene comentarios. Agrega
comentarios que:

\begin{enumerate}
\def\labelenumi{\arabic{enumi}.}
\tightlist
\item
  Expliquen qué hace cada sección
\item
  Documenten qué hace la función
\item
  Desactiven temporalmente una línea de código (sin eliminarla)
\end{enumerate}

\begin{Shaded}
\begin{Highlighting}[]
\KeywordTok{def}\NormalTok{ calcular\_total(precio, cantidad):}
    \ControlFlowTok{return}\NormalTok{ precio }\OperatorTok{*}\NormalTok{ cantidad}

\NormalTok{precio\_producto }\OperatorTok{=} \FloatTok{25.50}
\NormalTok{unidades }\OperatorTok{=} \DecValTok{4}
\NormalTok{total }\OperatorTok{=}\NormalTok{ calcular\_total(precio\_producto, unidades)}
\BuiltInTok{print}\NormalTok{(total)}
\BuiltInTok{print}\NormalTok{(}\StringTok{"Fin del cálculo"}\NormalTok{)}
\end{Highlighting}
\end{Shaded}

\begin{center}\rule{0.5\linewidth}{0.5pt}\end{center}

\textbf{Solución (no mires hasta haberlo intentado):}

\begin{Shaded}
\begin{Highlighting}[]
\CommentTok{\# Función que calcula el total de una compra}
\KeywordTok{def}\NormalTok{ calcular\_total(precio, cantidad):}
    \CommentTok{"""Multiplica el precio por la cantidad para obtener el total."""}
    \ControlFlowTok{return}\NormalTok{ precio }\OperatorTok{*}\NormalTok{ cantidad}

\CommentTok{\# Datos del producto}
\NormalTok{precio\_producto }\OperatorTok{=} \FloatTok{25.50}  \CommentTok{\# Precio unitario en euros}
\NormalTok{unidades }\OperatorTok{=} \DecValTok{4}             \CommentTok{\# Cantidad de unidades compradas}

\CommentTok{\# Cálculo del total}
\NormalTok{total }\OperatorTok{=}\NormalTok{ calcular\_total(precio\_producto, unidades)}

\CommentTok{\# Mostrar resultados}
\BuiltInTok{print}\NormalTok{(total)            }\CommentTok{\# Muestra el total}
\CommentTok{\# print("Fin del cálculo")  \# Desactivado temporalmente}
\end{Highlighting}
\end{Shaded}

\subsection{🗂️ Síntesis}\label{suxedntesis}

\begin{itemize}
\tightlist
\item
  \textbf{\texttt{\#}} inicia un comentario de una línea.
\item
  \textbf{Todo después de \texttt{\#}} es ignorado por Python.
\item
  \textbf{Comentarios al final de línea} explican código específico.
\item
  \textbf{Múltiples \texttt{\#}} para comentarios de varias líneas (o
  triple comilla).
\item
  \textbf{Desactivar código} comentando líneas es útil para pruebas.
\item
  \textbf{Docstrings} (\texttt{"""}) documentan funciones y clases.
\item
  \textbf{Buen comentario}: explica el ``por qué'', no repite el
  ``qué''.
\end{itemize}

\subsection{🔗 Qué sigue}\label{quuxe9-sigue}

\begin{itemize}
\tightlist
\item
  \textbf{Relacionado:} variables (cómo nombrar bien las variables para
  reducir necesidad de comentarios)
\item
  \textbf{Siguiente tema:} Variables --- aprenderás a almacenar datos
  con nombres claros y autoexplicativos.
\end{itemize}

\section{Publicaciones Similares}\label{publicaciones-similares}

Si te interesó este artículo, te recomendamos que explores otros blogs y
recursos relacionados que pueden ampliar tus conocimientos. Aquí te dejo
algunas sugerencias:

\begin{enumerate}
\def\labelenumi{\arabic{enumi}.}
\tightlist
\item
  \href{https://numerus-scriptum.netlify.app/python/2020-06-19-instalacion-de-anaconda/index.pdf}{\faIcon{file-pdf}}
  \href{https://numerus-scriptum.netlify.app/python/2020-06-19-instalacion-de-anaconda}{Instalacion
  De Anaconda}
\item
  \href{https://numerus-scriptum.netlify.app/python/2020-06-20-configurar-entorno-virtual-python-anaconda/index.pdf}{\faIcon{file-pdf}}
  \href{https://numerus-scriptum.netlify.app/python/2020-06-20-configurar-entorno-virtual-python-anaconda}{Configurar
  Entorno Virtual Python Anaconda}
\item
  \href{https://numerus-scriptum.netlify.app/python/2021-04-17-01-introducion-a-la-programacion-con-python/index.pdf}{\faIcon{file-pdf}}
  \href{https://numerus-scriptum.netlify.app/python/2021-04-17-01-introducion-a-la-programacion-con-python}{01
  Introducion A La Programacion Con Python}
\item
  \href{https://numerus-scriptum.netlify.app/python/2021-05-31-02-variables-expresiones-y-statements-con-python/index.pdf}{\faIcon{file-pdf}}
  \href{https://numerus-scriptum.netlify.app/python/2021-05-31-02-variables-expresiones-y-statements-con-python}{02
  Variables Expresiones Y Statements Con Python}
\item
  \href{https://numerus-scriptum.netlify.app/python/2021-06-07-03-objetos-de-python/index.pdf}{\faIcon{file-pdf}}
  \href{https://numerus-scriptum.netlify.app/python/2021-06-07-03-objetos-de-python}{03
  Objetos De Python}
\item
  \href{https://numerus-scriptum.netlify.app/python/2021-06-14-04-ejecucion-condicional-con-python/index.pdf}{\faIcon{file-pdf}}
  \href{https://numerus-scriptum.netlify.app/python/2021-06-14-04-ejecucion-condicional-con-python}{04
  Ejecucion Condicional Con Python}
\item
  \href{https://numerus-scriptum.netlify.app/python/2021-06-21-05-iteraciones-con-python/index.pdf}{\faIcon{file-pdf}}
  \href{https://numerus-scriptum.netlify.app/python/2021-06-21-05-iteraciones-con-python}{05
  Iteraciones Con Python}
\item
  \href{https://numerus-scriptum.netlify.app/python/2021-08-16-06-funciones-con-python/index.pdf}{\faIcon{file-pdf}}
  \href{https://numerus-scriptum.netlify.app/python/2021-08-16-06-funciones-con-python}{06
  Funciones Con Python}
\item
  \href{https://numerus-scriptum.netlify.app/python/2021-08-23-07-dataframes-con-python/index.pdf}{\faIcon{file-pdf}}
  \href{https://numerus-scriptum.netlify.app/python/2021-08-23-07-dataframes-con-python}{07
  Dataframes Con Python}
\item
  \href{https://numerus-scriptum.netlify.app/python/2021-11-29-08-prediccion-y-metrica-de-performance-con-python/index.pdf}{\faIcon{file-pdf}}
  \href{https://numerus-scriptum.netlify.app/python/2021-11-29-08-prediccion-y-metrica-de-performance-con-python}{08
  Prediccion Y Metrica De Performance Con Python}
\item
  \href{https://numerus-scriptum.netlify.app/python/2021-12-06-09-metodos-de-machine-learning-para-clasificacion-con-python/index.pdf}{\faIcon{file-pdf}}
  \href{https://numerus-scriptum.netlify.app/python/2021-12-06-09-metodos-de-machine-learning-para-clasificacion-con-python}{09
  Metodos De Machine Learning Para Clasificacion Con Python}
\item
  \href{https://numerus-scriptum.netlify.app/python/2021-12-13-10-metodos-de-machine-learning-para-regresion-con-python/index.pdf}{\faIcon{file-pdf}}
  \href{https://numerus-scriptum.netlify.app/python/2021-12-13-10-metodos-de-machine-learning-para-regresion-con-python}{10
  Metodos De Machine Learning Para Regresion Con Python}
\item
  \href{https://numerus-scriptum.netlify.app/python/2025-05-10-visualizacion-de-datos-con-python/index.pdf}{\faIcon{file-pdf}}
  \href{https://numerus-scriptum.netlify.app/python/2025-05-10-visualizacion-de-datos-con-python}{Visualizacion
  De Datos Con Python}
\end{enumerate}

Esperamos que encuentres estas publicaciones igualmente interesantes y
útiles. ¡Disfruta de la lectura!






\end{document}
